\begin{problem}[第35届全国中学生物理竞赛决赛][滑轮与软绳动力学]{98}
    如图，质量线密度为 $\lambda$ 、不可伸长的软细绳跨过一盘状定滑轮，定滑轮半径为 $R$ ，轴离地面高度为 $L$ 。系统原处于静止状态。在 $t = 0$ 时，滑轮开始以恒定角速度 $\omega$ 逆时针转动，绳子在滑轮带动下开始运动，绳子与滑轮间的动摩擦因数为 $\mu$ 。滑轮两侧的绳子在运动过程中始终可视为沿竖直方向，绳的两端在运动过程中均没有离开地面，地面上的绳子可视为集中在一点。已知重力加速度大小为 $g$ 。记绳子在与滑轮左、右侧相切处的张力大小分别为 $T_{1}$ 、 $T_{2}$ 。
    \begin{subproblem}
        分别列出在绳子速度达到最大值之前，滑轮两侧绳子的竖直部分及滑轮上任意一小段绳子的运动所满足的动力学方程组；
    \end{subproblem}
    \begin{subproblem}
        求绳子可达到的最大速度的大小。

        可参考的数学关系式:
        \begin{align*}
            & {\frac {\mathrm{d} y}{\mathrm{d} x} + \alpha   y = \mathrm{e} ^ {- \alpha x}   \frac {\mathrm{d} (y \mathrm{e} ^ {\alpha x})}{\mathrm{d} x};} \\
            & {\int \mathrm{e} ^ {\alpha x} \cos x \mathrm{d} x = \frac {\mathrm{e} ^ {\alpha x}}{1 + \alpha^ {2}} \big (\alpha \cos x + \sin x \big) + C _ {1}, \quad C _ {1} \text {为积分常数};} \\
            & {\int \mathrm{e} ^ {\alpha x} \sin x \mathrm{d} x = \frac {\mathrm{e} ^ {\alpha x}}{1 + \alpha^ {2}} \big (\alpha \sin x - \cos x \big) + C _ {2}, \quad C _ {2} \text {为积分常数。}}
        \end{align*}
    \end{subproblem}
\end{problem}